% BAN TIENG VIET cua 03_guarantee.tex (2026-10-08). So lieu giu nguyen.
\section{Khi nào được kết tội một nông dân?}
\label{sec:guarantee}

\subsection{Độ sâu chấp nhận được là một dải, không phải một con số}
\label{sec:umax}

Phiên bản đã công bố của lớp này nêu một giới hạn duy nhất, $u\le 61$~mm, là
khai báo lớn nhất không bao giờ gây báo động --- đó là sản phẩm của một hồ sơ
đất và một điều kiện đầu duy nhất. Với $V_4$ kích hoạt khi
$u_i>c_{t_i}+\varepsilon$ và $c_t=d_t\cdot\mathrm{span}$, giới hạn tại một
ruộng có độ cạn kiệt ban đầu $w_0$ và tưới không muộn hơn $d^{*}$ là
\begin{equation}
u_{\max}=\max(d^{*},w_0)\cdot\mathrm{span}+\varepsilon .
\label{eq:umax}
\end{equation}
Hai điều chỉnh rút ra, cả hai đều được xác minh bằng cách quét toán tử phát lại
chứ không chỉ bằng đại số (dẫn xuất ở Phụ lục~S14): \eqref{eq:umax} dùng
$\max$ chứ không dùng hiệu --- một ruộng khô ban đầu chấp nhận khai báo
\emph{lớn hơn} ruộng ướt, và biểu thức đã công bố trước đó không tương ứng với
trạng thái vật lý nào --- và nó chặn một \emph{lần tưới đơn lẻ} chứ không chặn
cả vụ, nên một đối thủ phát nhiều khai báo dưới ngưỡng có thể thổi phồng tổng
lượng mà không có giới hạn. Bảng~S2 của Phụ lục đưa ra dải phải được công bố
thay thế, $[49{,}0; 86{,}0]$~mm trên toàn $\Theta$.

\subsection{Kết tội bền vững dưới bất định mô hình}
\label{sec:robust}

Kiểm toán viên không biết đất của nông dân, không biết giai đoạn sinh trưởng,
cũng không biết hiệu suất của hệ thống dẫn nước. Gọi $\Theta$ là tập bất định
và $V_\theta(L)$ là tập vi phạm dưới cấu hình $\theta$; bộ phát hiện đã công bố
chỉ đánh giá một $\theta_0$ mặc định và coi mọi vi phạm là có tội. Chúng tôi
thay nó bằng một phép giao:
\begin{equation}
\begin{aligned}
V_{\mathrm{rob}}(L)&=\bigcap_{\theta\in\Theta}V_\theta(L),\\
\Theta&=\mathcal{S}\times\mathcal{K}\times\mathcal{E}\times\mathcal{W},
\end{aligned}
\label{eq:robust}
\end{equation}
với $|\mathcal{S}|=5$ loại đất, $|\mathcal{K}|=2$ chế độ hệ số cây trồng,
$|\mathcal{E}|=3$ mức hiệu suất tưới và $|\mathcal{W}|=5$ trạng thái đầu, nên
$|\Theta|=150$ và mỗi nhật ký tốn $150$ lần phát lại. Ngữ nghĩa được chọn cố ý
lệch về một phía:
\begin{equation}
a(L)=
\begin{cases}
\text{kết tội}, & V_{\mathrm{rob}}(L)\neq\emptyset,\\
\text{chờ xác minh}, & \text{trường hợp còn lại}.
\end{cases}
\label{eq:verdict}
\end{equation}

\begin{proposition}[Không kết tội oan trong $\Theta$]
\label{prop:noalarm}
Nếu cấu hình thật của ruộng $\theta^{*}$ thuộc $\Theta$ và $L$ tuân thủ dưới
$\theta^{*}$, thì $V_{\mathrm{rob}}(L)=\emptyset$ và nông dân không bị kết tội.
\end{proposition}

\begin{proof}
$V_{\mathrm{rob}}(L)\subseteq V_{\theta^{*}}(L)=\emptyset$.
\end{proof}

Mệnh đề~\ref{prop:noalarm} là thứ làm cho lớp này triển khai được: sai số mô
hình có thể làm giảm năng lực phát hiện nhưng không thể tạo ra một lời kết
tội, với một cái giá phải được nói rõ. Hai chi tiết cài đặt --- cách xử lý các
thế giới không kích hoạt và độ trải của $\mathcal{W}$ --- sẽ âm thầm phá vỡ
mệnh đề nếu làm sai (Phụ lục~S7).

\subsection{Ngưỡng chênh lệch bơm tỉ lệ thuận}

Quy tắc tuyệt đối $G>\tau_0$ với $\tau_0=5$~mm giả định nhật ký là bản sao
trung thực của lượng nước đã bơm. Nó không phải vậy: quên ghi, làm tròn đến
$5$~mm gần nhất và ghi trễ một ngày đều đóng góp vào $G$ mà không có hành vi
sai trái nào, nên ngưỡng phải co giãn theo số lần tưới được ghi:
\begin{equation}
\begin{aligned}
\tau(n)&=\delta_{\mathrm{rel}}M+z\sqrt{n}\,\sigma_{\mathrm{event}},\\
\sigma_{\mathrm{event}}^{2}&=p(1-p)\bar u^{2}+\tfrac{r^{2}}{12}
+q(1-q)(s\bar u)^{2},
\end{aligned}
\label{eq:tau}
\end{equation}
với $p$ là tỷ lệ quên, $r$ là bước làm tròn, $q$ là tỷ lệ sai độ sâu và $s$ là
biên độ của nó, $\bar u=30$~mm là độ sâu điển hình của một lần tưới,
$\delta_{\mathrm{rel}}=2\%$ và $z=3$; quy tắc tuyệt đối là trường hợp suy biến
$n\to0$, $\sigma_{\mathrm{event}}\to0$.

\begin{remark}[Một bộ lọc không thể kích hoạt]
\label{sec:v2reach}
Bộ lọc vật lý thứ ba của phiên bản đã công bố, đánh dấu độ cạn kiệt xuống dưới
điểm héo, đã không hoạt động theo hai cách: ngưỡng của nó ($1{,}02$) nằm ngoài
dải đã cắt ngọn của chỉ số được báo cáo ($0$ lần kích hoạt trong $281$ nhật
ký), và sau khi sửa thì nó vẫn không thể kích hoạt --- qua $150$ cấu hình với
một nhật ký rỗng, độ cạn kiệt chưa cắt ngọn chỉ đạt đỉnh $0{,}9987$, và một
ruộng lúa vụ khô ở đồng bằng không chạm tới điểm héo (Bảng~S3 của Phụ lục, dẫn
xuất ở S13). Lớp này có hai bộ lọc vật lý hoạt động và một thiết bị an toàn
bất động.
\end{remark}
